\documentclass[10pt,oneside]{book}
\newcommand{\gtitle}{Theory of Computation - Week 3}
%%%% Page Info + Commands %%%%%{

%packages
\usepackage{pgfplots}
\pgfplotsset{compat=1.18}
\usepackage{amsmath}
\usepackage{amsfonts}
\usepackage{charter}
\usepackage{amsthm,letltxmacro}
\usepackage{amssymb}
\usepackage{fancyhdr}
\usepackage{float}
\usepackage[most]{tcolorbox}
\usepackage{enumerate}
\usepackage{enumitem}
\usepackage[dvipsnames]{xcolor}
\usepackage{changepage}
\usepackage{mdframed}
\usepackage{graphicx}

% SMALL PAGE COMMAND
\newcommand{\smallpage}{  \setlength \oddsidemargin{.25in}
            \setlength \textwidth{5.5in}
            \setlength \topmargin{0in}
            \setlength \textheight{9in}}


% Commands for sets of numbers
\newcommand{\Z}{\mathbb{Z}}\newcommand{\R}{\mathbb{R}}\newcommand{\N}{\mathbb{N}}\newcommand{\Q}{\mathbb{Q}}\newcommand{\C}{\mathbb{C}}
\newcommand{\real}[1]{\text{Re}\left(#1\right)}
\newcommand{\im}[1]{\text{Im}\left(#1\right)}

% gordie spacing and boxing commands
\newcommand{\gspc}{\vspace{0.13cm}} % 0.5 space
\newcommand{\gline}[0]{\gspc\\} % 1.5 space
\newcommand{\gbline}{\gspc\gspc\\} % double space
\newcommand{\gbox}[1]{\fbox{\parbox{\linewidth}{#1}}} % multiline box command
\newcommand{\gmod}[1]{\,(\text{mod}\,#1)}


\newcommand{\gimage}[2]{
\begin{figure}[H]
    \centering
    \includegraphics[width=#2]{#1}
\end{figure}
}

\newcommand{\centerit}[1]{{\begin{center} #1 \end{center}}}
\newcommand{\ul}[1]{\underline{#1}}


% COLOR BOX
\newmdenv[backgroundcolor=gray!8,
  linewidth=0pt, innerleftmargin=0pt, innerrightmargin=0pt,
  skipabove=\baselineskip, skipbelow=\baselineskip
]{coloredadjust}

% PROOF
\LetLtxMacro\oldproof\proof
\let\endoldproof\endproof
\renewenvironment{proof}[1][\proofname]
  {\vspace{0.2cm}\begin{adjustwidth}{2em}{2em}
   \oldproof[#1]}
  {\endoldproof
\end{adjustwidth}}

% PROBLEM
\newenvironment{problem}[1]
  {\vspace{-0.1cm}\begin{coloredadjust}\begin{adjustwidth}{2em}{2em}
   \textit{Problem. }#1}
  {\endoldproof
\end{adjustwidth}\end{coloredadjust}\gspc}

%% DEFINITION
\newtcbtheorem[]{definition}{Definition}{enhanced,
  colback=blue!2, colframe=blue!50!black, coltitle=blue!70!black,
  fonttitle=\bfseries,
  boxrule=0pt, toprule=2pt, bottomrule=1pt, leftrule=1pt, rightrule=1pt, arc=1pt, outer arc=1pt,
  attach title to upper, title={#1},
}{dfn}

%% CRITERION
\newtcbtheorem[]{criterion}{Criterion}{enhanced,
  colback=magenta!3, colframe=magenta!60!black, coltitle=magenta!20!black,
  fonttitle=\bfseries,
  boxrule=4pt, toprule=2pt, bottomrule=1pt, leftrule=1pt, rightrule=1pt, arc=1pt, outer arc=1pt,
  attach title to upper, title={#1},
}{ctn}

%% COROLLARY
\newtcbtheorem[]{corollary}{Corollary}{enhanced,
  colback=orange!3, colframe=orange!80!black, coltitle=orange!50!black,
  fonttitle=\bfseries,
  boxrule=4pt, toprule=1pt, bottomrule=1pt, leftrule=1pt, rightrule=1pt, arc=5pt, outer arc=5pt,
  attach title to upper, title={#1},
}{ctn}

%% THEOREM
\newtcbtheorem[]{theorem}{Theorem}{enhanced,
  colback=green!2, colframe=green!40!black, coltitle=green!20!black,
  fonttitle=\bfseries,
  boxrule=4pt, toprule=3pt, bottomrule=1pt, leftrule=1pt, rightrule=1pt,arc=5pt, outer arc=5pt,
  attach title to upper, title={#1},
}{thm}

%%%% Page Setup %%%%%%%
\smallpage\sloppy
\pagestyle{fancy}
\lhead{\large{\textbf{\gtitle}}}
%\rfoot{\thepage/\pageref{LastPage} }
\setlength{\headheight}{10pt}

\begin{document}

\newcommand{\cg}[1]{\color{OliveGreen}#1\color{white}}
\newcommand{\cb}[1]{\color{BlueViolet}#1\color{white}}

\subsection*{Warm Up}
Let $A$ be the set of binary string that contains at least two 0s.

Let $B$ be the set of binary strings that end in 01. 

\begin{enumerate}
  \item Draw a diagram for a DFA that recognizes $A^*$. 
  \gimage{image.png}{5cm}
  \item Draw a diagram for a DFA that recognizes $B^*$
  \gimage{image2.png}{5cm}
\end{enumerate}
\subsection*{Hypothesis}
\textbf{Prediction: } If $A$ and $B$ are regular languages, then:
\begin{align*}
  A\cup B,\;A\circ B\text{ and } A^*
\end{align*}  
are regular languages as well. \gline{}
In order to prove this, is that we must construct a DFA that proves both. 
\subsection*{Sketch of Proof}
Let $A$ be the set of binary string that contain at least two 0s. Let $B$ be the set of binary strings that end in 01. Is this a DFA for $A\cup B$. 
\gimage{image3.png}{6cm}
\subsection*{Nondeterminism}
In a \textbf{deterministic} finit automaton, we have exactly one way to proceed from each state for each character in the alphabet.

In a \textbf{nondeterministic} finite automaton, there can only be \textbf{one, many, or no} ways to proceed from each state for each character in the alphabet. 
\subsection*{Exercise}
What language is recognized by the NFA pictured below?
\gimage{image4.png}{5cm}
It's strings of the form 00, 010, or 01$1^*$.
\subsection*{$\varepsilon$-Arrows}
Nondeterministic finite automata also allow us to ahve $\varepsilon$-arrows, where we proceed without reading a character.

\begin{definition}{NFA - Formal Definition}{}\gline{}
  A \textbf{nondeterministic finite automaton} is a 5-tuple $(Q, \sum, \delta, q_0, F)$, where:
  \begin{itemize}
    \item $Q$ is a finite set called the \textbf{states},
    \item $\sum$ is a finite set called the \textbf{alphabet}
    \item $\delta:\;Q\times \sum_\varepsilon\to P(Q)$ is the \textbf{transition function}
    \item $q_0$ in $Q$ is the \textbf{start state}, and,
    \item $F\subseteq Q$ is the set of \textbf{accept states}
  \end{itemize}
  Here, $\sum_\varepsilon\to P(Q)$ denotes $\sum U\{\varepsilon\}$ (the alphabet, plus $\varepsilon$). 

  $P(Q)$ is the power set of $Q$ (the set of all subsets of $Q$)
\end{definition}

\section*{Last time in Theory of Computation}
Is the following automata an NFA or a DFA, why?

\end{document}